\documentclass[a4,11pt]{aleph-notas}

% -- Paquetes adicionales
\usepackage{enumitem}
\usepackage{aleph-comandos}

% -- Datos
\institucion{Facultad de Ciencias Exactas, Naturales y Ambientales}
\carrera{Catálogo STEM}
\asignatura{Matemática III}
\tema{Resumen 09: Rotacional, Divergencia y Campos conservativos}
\autor[A. Merino]{Andrés Merino}
\fecha{Periodo 2026-1}

\logouno[0.16\textwidth]{Logos/logoPUCE_04_ac}
\definecolor{colortext}{HTML}{1957d1}
\definecolor{colordef}{HTML}{1957d1}
\fuente{montserrat}


\begin{document}

\encabezado

En este resumen consideraremos $\Omega\subset \R^n$ un conjunto abierto.

%%%%%%%%%%%%%%%%%%%%%%%%%%%%%%%%%%%%%%
\section{Divergencia y rotacional}
%%%%%%%%%%%%%%%%%%%%%%%%%%%%%%%%%%%%%%

\begin{defi}[Divergencia]
    Dado un campo vectorial $\func{F}{\Omega}{\R^n}$ diferenciable, se define la divergencia de $F$ por
    \[
        \dive(F)(x) = \sum_{i=1}^n D_i f_i(x) = \sum_{i=1}^n \parcial{f_i}{x_i}(x)
    \]
    para todo $x\in\Omega$.
\end{defi}

\begin{defi}
    Dados un campo vectorial $\func{F}{\Omega}{\R^n}$ diferenciable y $x\in\Omega$:
    \begin{itemize}
        \item si $\dive(F)(x)>0$, se dice que $x$ es una fuente;
        \item si $\dive(F)(x)<0$, se dice que $x$ es un sumidero.
    \end{itemize}
    Si $\dive(F)=0$ en todo $\Omega$, se dice que $F$ es un campo incompresible.
\end{defi}

\begin{defi}[Rotacional en $\R^2$]
    Si $n=2$, dado un campo vectorial $\func{F}{\Omega\subset \R^2}{\R^2}$ diferenciable, se define el rotacional de $F$ por
    \[
        \rot(F)(x,y) = D_1 f_2(x,y) - D_2 f_1(x,y) = \parcial{f_2}{x}(x,y) - \parcial{f_1}{y}(x,y)
    \]
    para $(x,y)\in \Omega$. Si $\rot(F)(x,y)>0$ el campo tiene rotación antihoraria en ese punto; si $\rot(F)(x,y)<0$, rotación horaria.
\end{defi}

\begin{defi}[Rotacional en $\R^3$]
    Si $n=3$, dado un campo vectorial $\func{F}{\Omega}{\R^3}$ diferenciable, se define el rotacional de $F$ por
    \begin{align*}
        \rot(F)(x,y,z)
        & = \begin{pmatrix}
                \parcial{f_3}{y}(x,y,z) - \parcial{f_2}{z}(x,y,z) \\[4mm]
                \parcial{f_1}{z}(x,y,z) - \parcial{f_3}{x}(x,y,z) \\[4mm]
                \parcial{f_2}{x}(x,y,z) - \parcial{f_1}{y}(x,y,z)
            \end{pmatrix}^T
    \end{align*}
    para $(x,y,z)\in \Omega$. Si $\rot(F)=0$ en todo $\Omega$, se dice que $F$ es un campo irrotacional.
\end{defi}

\begin{advertencia}
    Se suele usar la notación
    \[
        \dive(F) = \nabla\cdot F
        \qquad\text{y}\qquad
        \rot(F) = \nabla \times F.
    \]
\end{advertencia}

\begin{prop}
    Si $n=3$, dado $\func{f}{\Omega\subset\R^3}{\R}$ un campo escalar diferenciable, se tiene que
    \[
        \rot(\nabla f)(x,y,z)=0
    \]
    para $(x,y,z)\in \Omega$.
\end{prop}

\begin{prop}
    Si $n=3$, dado $\func{F}{\Omega\subset\R^3}{\R^3}$ un campo vectorial diferenciable, se tiene que
    \[
        \dive(\rot(F))(x,y,z)=0
    \]
    para $(x,y,z)\in \Omega$.
\end{prop}

\begin{defi}[Laplaciano]
    Dado un campo escalar $\func{f}{\Omega}{\R}$ diferenciable, se define el laplaciano de $f$ por
    \[
        \Delta f (x) = \dive(\nabla f)(x) = \sum_{i=1}^n \frac{\partial^2 f}{\partial x_i^2}(x)
    \]
    para todo $x\in\R^n$. Se suele denotar $\Delta f = \nabla^2 f$.
\end{defi}

%%%%%%%%%%%%%%%%%%%%%%%%%%%%%%%%%%%%%%
\section{Campos conservativos}
%%%%%%%%%%%%%%%%%%%%%%%%%%%%%%%%%%%%%%

\begin{defi}
    Dado un campo vectorial $\func{F}{\Omega}{\R^n}$, se dice que es conservativo si existe un campo escalar $\func{f}{\Omega}{\R}$ tal que
    \[
        F(x)=\nabla f(x)
    \]
    para todo $x\in \Omega$. De ser este el caso, se dice que $f$ es un potencial de $F$.
\end{defi}

\begin{prop}
    Sea $\func{F}{\Omega}{\R^n}$ un campo vectorial diferenciable. Si $F$ es conservativo, entonces
    \[
        D_i F_j(x) = D_j F_i(x)
    \]
    para todo $x\in \Omega$ y todo $i,j=1,\ldots,n$.
\end{prop}

\begin{defi}
    Dado $\Omega\subset \R^n$, se dice que $\Omega$ es simplemente conexo si toda curva cerrada puede ser deformada continuamente a un punto en $\Omega$. La idea intuitiva es que $\Omega$ no posee ``huecos''.
\end{defi}

\begin{prop}
    Sean $\func{F}{\Omega}{\R^n}$, con $n=2,3$, un campo vectorial diferenciable donde $\Omega$ es simplemente conexo. Entonces $F$ es conservativo si y solo si
    \[
        D_i F_j(x) = D_j F_i(x)
    \]
    para todo $x\in \Omega$ y todo $i,j=1,\ldots,n$.
\end{prop}

\begin{prop}
    Si $n\in\{2,3\}$ y $\func{F}{\Omega}{\R^n}$ es un campo vectorial con $\rot(F)=0$ y $\Omega$ simplemente conexo, entonces $F$ es conservativo.
\end{prop}


\end{document}
